\begin{table}[H]\centering\begin{tabular}{>{\centering\arraybackslash}m{1.5cm}|>{\centering\arraybackslash}m{1.2cm}>{\centering\arraybackslash}m{1.2cm}|>{\centering\arraybackslash}m{1.2cm}>{\centering\arraybackslash}m{1.2cm}|>{\centering\arraybackslash}m{1.2cm}>{\centering\arraybackslash}m{1.2cm}}
\toprule
elements & \multicolumn{2}{|c}{gravity} & \multicolumn{2}{|c}{traction} & \multicolumn{2}{|c}{compression} \\
 & B & G & B & G & B & G \\
\midrule
24,576 & 0.074 & 0.055 & 0.055 & 0.100 & 0.072 & 0.084 \\
63,888 & 0.089 & 0.095 & 0.086 & 0.133 & 0.083 & 0.093 \\
162,000 & 0.105 & -0.030 & 0.119 & 0.194 & 0.029 & 0.085 \\
444,528 & -0.100 & 0.224 & 0.072 & 0.066 & 0.082 & 0.123 \\
\bottomrule
\end{tabular}\caption{Estimated first-run overhead (in seconds) for different load types and mesh sizes using \textbf{PyTorch} under \textbf{FP32} precision. The overhead is computed as the difference between the first execution time and the average steady-state runtime over the subsequent 7 runs. This metric captures framework-specific initialization costs, including compilation (where applicable), kernel setup, memory allocation, and cache warm-up effects.}\label{tab:pytorch_overhead_fp32}\end{table}
\begin{table}[H]\centering\begin{tabular}{>{\centering\arraybackslash}m{1.5cm}|>{\centering\arraybackslash}m{1.2cm}>{\centering\arraybackslash}m{1.2cm}|>{\centering\arraybackslash}m{1.2cm}>{\centering\arraybackslash}m{1.2cm}|>{\centering\arraybackslash}m{1.2cm}>{\centering\arraybackslash}m{1.2cm}}
\toprule
elements & \multicolumn{2}{|c}{gravity} & \multicolumn{2}{|c}{traction} & \multicolumn{2}{|c}{compression} \\
 & B & G & B & G & B & G \\
\midrule
24,576 & 0.370 & 0.412 & 0.483 & 0.545 & 0.499 & 0.514 \\
63,888 & 0.427 & 0.401 & 0.568 & 0.526 & 0.554 & 0.515 \\
162,000 & 0.394 & 0.406 & 0.519 & 0.518 & 0.528 & 0.526 \\
444,528 & 0.435 & 0.415 & 0.485 & 0.524 & 0.570 & 0.542 \\
\bottomrule
\end{tabular}\caption{Estimated first-run overhead (in seconds) for different load types and mesh sizes using \textbf{JAX} under \textbf{FP32} precision. The overhead is computed as the difference between the first execution time and the average steady-state runtime over the subsequent 7 runs. This metric captures framework-specific initialization costs, including compilation (where applicable), kernel setup, memory allocation, and cache warm-up effects.}\label{tab:jax_overhead_fp32}\end{table}
\begin{table}[H]\centering\begin{tabular}{>{\centering\arraybackslash}m{1.5cm}|>{\centering\arraybackslash}m{1.2cm}>{\centering\arraybackslash}m{1.2cm}|>{\centering\arraybackslash}m{1.2cm}>{\centering\arraybackslash}m{1.2cm}|>{\centering\arraybackslash}m{1.2cm}>{\centering\arraybackslash}m{1.2cm}}
\toprule
elements & \multicolumn{2}{|c}{gravity} & \multicolumn{2}{|c}{traction} & \multicolumn{2}{|c}{compression} \\
 & B & G & B & G & B & G \\
\midrule
24,576 & 0.021 & 0.004 & 0.007 & 0.005 & 0.003 & 0.010 \\
63,888 & 0.018 & 0.030 & -0.021 & -0.001 & 0.007 & 0.020 \\
162,000 & -0.011 & 0.051 & -0.025 & -0.015 & 0.013 & 0.028 \\
444,528 & 0.276 & -0.032 & -0.047 & 0.144 & 0.056 & -0.005 \\
\bottomrule
\end{tabular}\caption{Estimated first-run overhead (in seconds) for different load types and mesh sizes using \textbf{Warp} under \textbf{FP32} precision. The overhead is computed as the difference between the first execution time and the average steady-state runtime over the subsequent 7 runs. This metric captures framework-specific initialization costs, including compilation (where applicable), kernel setup, memory allocation, and cache warm-up effects.}\label{tab:warp_overhead_fp32}\end{table}
